\documentclass{article}
\setlength{\topmargin}{-0.25 in}
\setlength{\textwidth}{6.5in}
\setlength{\textheight}{8.5 in}
\setlength{\oddsidemargin}{0in}
\setlength{\evensidemargin}{0in}
\setlength{\parindent}{0.0in}
\setlength{\parskip}{7pt plus 1pt}

\usepackage{textcomp}
\usepackage{bcprules}
\usepackage{xcolor}				% Allow for color (annotations)
\usepackage{stmaryrd}
\usepackage{amssymb}
\usepackage{hyperref}

\def\sectionautorefname{Section}
\def\subsectionautorefname{Section}
\def\subsubsectionautorefname{Section}
\def\tableautorefname{Table}
\def\figureautorefname{Figure}

\def\tdm#1{\textcolor{blue}{\textsf{TDM: #1}}}
\def\ab#1{\textcolor{red}{\textsf{AB: #1}}}
\def\tf#1{\textcolor{purple}{\textsf{TF: #1}}}

\newcommand{\maybe}{{\bf\footnotesize{?}}\kern-1pt}

\newcommand{\eval}[3]{{\tt #1} $\vdash$ {\tt #2} $\Downarrow$ {\tt #3}}
\newcommand{\stay}[2]{stay({\tt #1}) $=$ {\tt #2}}
\newcommand{\toSolver}[2]{{\tt E} $\vdash$ {\tt #1} $\leadsto$ $#2$}
\newcommand{\fromSolver}[2]{$#1$ $\leadsto$ {\tt #2}}
\newcommand{\csolve}[2]{$#1$ $\models$ $#2$}
\newcommand{\solve}[2]{{\tt #1} $\models$ {\tt #2}}
\newcommand{\cfail}[1]{$\not\models$ $#1$}
\newcommand{\fail}[1]{$\not\models$ {\tt #1}}
\newcommand{\step}[6]{{\tt <#1|#2|#3>} $\longrightarrow$
  {\tt <#4|#5|#6>}}
\newcommand{\tc}[3]{{\tt #1} $\vdash$ {\tt #2} : {\tt #3}}
\newcommand{\tcC}[2]{{\tt #1} $\vdash$ {\tt #2}}

\newcommand{\evalH}[4]{{\tt #1;#2} $\vdash$ {\tt #3} $\Downarrow$ {\tt #4}}
\newcommand{\solveH}[3]{{\tt #1;#2} $\models_{E,H}$ {\tt #3}}
\newcommand{\failH}[1]{$\not\models_{E,H}$ {\tt #1}}
\newcommand{\stepH}[8]{{\tt <#1|#2|#3|#4>} $\longrightarrow$
  {\tt <#5|#6|#7|#8>}}
\newcommand{\tcH}[4]{{\tt #1;#2} $\vdash$ {\tt #3} : {\tt #4}}
\newcommand{\tcHC}[3]{{\tt #1;#2} $\vdash$ {\tt #3}}
\newcommand{\toSolverH}[2]{{\tt E,H} $\vdash$ {\tt #1} $\leadsto$ $#2$}
\newcommand{\fromSolverH}[3]{$#1$ $\leadsto$ {\tt #2,#3}}

\newcommand{\toSolverOO}[6]{{\tt #1,#2} $\vdash$ {\tt #3} $\leadsto$ {\tt #4,#5} $\mid$ $#6$}

\newcommand{\inferrule}[3]{\infrule[#1]{\mbox{#2}}{\mbox{#3}}}
\newcommand{\inferax}[2]{\infax[#1]{\mbox{#2}}}

\newcommand{\judgment}[1]{\begin{flushleft} \fbox{#1}\end{flushleft}}

\newcommand{\receq}{\stackrel{{\tt r}}{{\tt =}}}

\sloppy

\begin{document}

\title{New Semantics for Babelsberg}
\author{Tim Felgentreff, Todd Millstein, and Alan Borning}
\date{}

% changed title to make it different from the title of the
% hypothetical POPL submission

\maketitle

\section{Introduction}

In this semantics the goal is to understand the restrictions we'd like
to put on constraint solving that arise solely due to the desire to
limit surprises for programmers.  To that end, in this semantics we
assume that the constraint solver directly ``understands'' the
constraints that can be expressed at the source level, with no need
for any kind of translation.


\section{Babelsberg/PrimitiveTypes}
\label{sec:babelsberg/int}


We present the formal semantics of Babelsberg/PrimitiveTypes. Besides
the program, there is an \emph{environment} that maps names to values.

\subsection{Syntax}
\label{sec:babelsberg-primitivetypes-syntax}

\begin{tt}
\begin{tabular}{lrrl}
{\rm Statement} & s & ::= & skip $\mid$ x := e $\mid$ always C $\mid$ once C $\mid$ s;s $\mid$ fail \\
  & & & $\mid$ if e then s else s $\mid$ while e do s \\
{\rm Constraint} & C & ::= & e $\mid$ $\rho$ e \\
{\rm Expression} & e & ::= & c $\mid$ x $\mid$ x? $\mid$ e $\oplus$ e  $\mid$ e $\olessthan$ e $\mid$ e $\ovee$ e\\
{\rm Constant} & c & ::= & true $\mid$ false $\mid$ nil $\mid$ {\rm base type constants}\\
{\rm Variable} & x & ::= & {\rm variable names} \\\\

{\rm Value} & v & ::= & c
\end{tabular}
\end{tt}

The language includes a set of boolean and real constants and the
undefined value {\tt nil}, ranged over by metavariable {\tt c}. A
finite set of operators on expressions is ranged over by $\oplus$. For
booleans this set includes no operations, but it does include
operations on the reals such as {\tt +}, {\tt *}, and {\tt   \textsurd}. 
The symbol $\olessthan$ ranges over a set of {\em predicate} operators
({\tt $=$} and {\tt $\neq$} for booleans, $\leq$, $<$, $=$, and so on
for reals.)  These test properties of base types in the language. The
symbol $\ovee$ ranges over a set of logical operators for combining
boolean expressions (e.g., {\tt $\wedge$}, {\tt $\vee$}).  The
predicate operators are assumed to include at least an equality
operator {\tt =} for each primitive type in the language, and the
logical operators are assumed to include at least conjunction {\tt $\wedge$}.
The syntax of this language does have some limitations as compared
with that of a practical language --- for example, there are only
binary operators (not unary or ternary), and the result must have the
same type as the arguments.  We make these simplifications since the
purpose of Babelsberg/PrimitiveTypes is to elucidate the semantics of
such languages as as a step toward Babelsberg/Objects, rather than to
specify a real language.

For constraints, the symbol $\rho$ ranges over a finite and totally
ordered set of constraint {\em priorities} and is assumed to include a
bottom element {\em weak}, an intermediate element {\em strong} and
top element {\em required}.  A constraint with no explicit priority is
assumed to have the priority {\em required}.

The syntax is thus that of a simple, 
standard imperative language except for the
{\tt always} and {\tt once} statements, which declare constraints, and
the {\tt ?} annotation for variables. An {\tt always} constraint must
hold for the rest of the programs execution, whereas a {\tt once}
constraint is satisfied by the solver and then retracted. The {\tt ?}
indicates that a variable is read-only if used in a constraint,
i.e.\@, the programmer does not allow the solver to change the
variables in the expression to satisfy constraints. Finally, the {\tt
  fail} statement is not intended for use at the source level, but
rather is used in the semantics to make explicit that the program
halts due either to an unsatisfiable constraint or to the solver being
unable to determine whether the constraint is satisfiable.  (The
formal semantics doesn't distinguish between these two cases, but
lumps them both together as {\tt fail}.  In a practical
implementation, we would likely want to distinguish them, since it's
useful if we can inform the programmer that the constraints are truly
not satisfiable.  We could also add standard exception handling to
remove the unsatisfiable or unknown constraint and continue, but omit
this here for simplicity.)

\subsection{Semantics}

The semantics is defined by several judgments, defined below.  These
judgments depend on the notion of an {\em environment}, which is a
partial function from program variables to program
values. Metavariable {\tt E} ranges over environments. We use {\tt
  E(x)$\uparrow$} to indicate that {\tt E} has no mapping for {\tt
  x}. When convenient we also view an environment as a set of (program
variable, program value) pairs.  For each operator {\tt o} in the language
we assume the existence of a corresponding semantic function denoted
$\llbracket {\tt o} \rrbracket$.

\judgment{\eval{E}{e}{v}}

``Expression {\tt e} evaluates to value {\tt v} in the context of environment {\tt
  E}.''

\inferax{E-Const}{\eval{E}{c}{c}}
\inferrule{E-Var}{
  {\tt E(x) = v}
}{
  \eval{E}{x}{v}
}


\inferrule{E-ReadOnly}{
  {\tt E(x) = v}
}{
  \eval{E}{x?}{v}
}

\inferrule{E-Op}{
  \eval{E}{e$_1$}{v$_1$} \andalso
  \eval{E}{e$_2$}{v$_2$} \andalso
  {\tt v$_1$ $\llbracket \oplus \rrbracket$ v$_2$} = {\tt v}
}{
  \eval{E}{e$_1$ $\oplus$ e$_2$}{v}
}

\inferrule{E-Compare}{
  \eval{E}{e$_1$}{v$_1$} \andalso
  \eval{E}{e$_2$}{v$_2$} \andalso
  {\tt v$_1$ $\llbracket \olessthan \rrbracket$ v$_2$} = {\tt v}
}{
  \eval{E}{e$_1$ $\olessthan$ e$_2$}{v}
}

\inferrule{E-Combine}{
  \eval{E}{e$_1$}{v$_1$} \andalso
  \eval{E}{e$_2$}{v$_2$} \andalso
  {\tt v$_1$ $\llbracket \ovee \rrbracket$ v$_2$} = {\tt v}
}{
  \eval{E}{e$_1$ $\ovee$ e$_2$}{v}
}


\judgment{\solve{E}{C}}

This judgment represents a call to the constraint solver, which we
treat as a black box.  The proposition 
\solve{E}{C} denotes that environment {\tt E} is a {\em solution} to the constraint
{\tt C} (and further one that is optimal according to the solver's
semantics, as discussed earlier).  

We use the notation \fail{C} to mean that there is no environment {\tt E}
such that \solve{E}{C}, i.e., the solver is unable to find a
solution, either because {\tt C} is unsatisfiable or because it is
satisfiable but the solver is unable to find one.

\judgment{\stay{E}{C}}

This judgment defines how to translate an environment into a
source-level ``stay'' constraint.

\inferax{StayEmpty}{\stay{$\emptyset$}{true}}

\inferrule{StayOne}{${\tt E}({\tt x}) = {\tt v}$ \andalso ${\tt E}_0 =
  {\tt E} \setminus \{({\tt x},{\tt v})\}$
  \andalso \stay{E$_0$}{C$_0$} \andalso {\tt C} = {\tt C$_0$}
  $\wedge$ {\tt weak x=v}}
{\stay{E}{C}}

We provide judgments to support the notion of different types for
expressions. These provide a non-standard form of type-checking to
check whether a constraint expression is valid (i.e., it returns a
boolean.) It is non-standard because it only happens dynamically on a
constraint, in the context of the current run-time environment.  The
types are also nonstandard: we only distinguish between booleans,
undefined values, and lump all other primitive types together.


\judgment{\step{E}{C}{s}{E$'$}{C$'$}{s$'$}}

``Configuration {\tt <E|C|s>} takes one step of execution to state {\tt <E$'$|C$'$|s$'$>}.''

A ``configuration'' defining the state of an execution includes a
concrete context, represented by the environment, a symbolic context,
represented by the constraint, and a statement to be executed.  The
environment and statement are standard, while the constraint is not
part of the state of a computation in most languages.  Intuitively,
the environment comes from constraint solving during the evaluation of
the immediately preceding statement, and the constraint records the
\verb|always| constraints that have been declared so far during
execution. 

\inferrule{S-Asgn}{
  \eval{E}{e}{v} \andalso
  \stay{E}{C$_s$} \andalso
  \solve{E$'$}{(C $\wedge$ C$_s$ $\wedge$ x = v)}
}{
  \step{E}{C}{x := e}{E$'$}{C}{skip}
}

\inferrule{S-AsgnFail}{
  \eval{E}{e}{v} \andalso
  \stay{E}{C$_s$} \andalso
  \fail{(C $\wedge$ C$_s$ $\wedge$ x = v)}
}{
  \step{E}{C}{x := e}{E}{C}{fail}
}

\inferrule{S-Once}{
  \stay{E}{C$_s$} \andalso
  \solve{E$'$}{(C $\wedge$ C$_s$ $\wedge$ C$_0$)}
}
{\step{E}{C}{once C$_0$}{E$'$}{C}{skip}}

\inferrule{S-OnceFail}{
  \stay{E}{C$_s$} \andalso
  \fail{(C $\wedge$ C$_s$ $\wedge$ C$_0$)}
}{
  \step{E}{C}{once C$_0$}{E}{C}{fail}
}

\inferrule{S-Always}{
  {\tt C$'$} = {\tt (C $\wedge$ C$_0$)} \andalso \stay{E}{C$_s$} \andalso
  \solve{E$'$}{(C$'$ $\wedge$ C$_s$)}
}{
  \step{E}{C}{always C$_0$}{E$'$}{C$'$}{skip}
}

\inferrule{S-AlwaysFail}{
  {\tt C$'$} = {\tt (C $\wedge$ C$_0$)} \andalso \stay{E}{C$_s$} \andalso
  \fail{(C$'$ $\wedge$ C$_s$)}
}{
  \step{E}{C}{always C$_0$}{E}{C}{fail}
}

\inferax{S-SeqSkip}{\step{E}{C}{skip;s$_2$}{E}{C}{s$_2$}}

\inferrule{S-SeqStep}{
  \step{E}{C}{s$_1$}{E$'$}{C$'$}{s$_1'$}
}{
  \step{E}{C}{s$_1$;s$_2$}{E$'$}{C$'$}{s$_1'$;s$_2$}
}

\inferax{S-SeqFail}{\step{E}{C}{fail;s$_2$}{E}{C}{fail}}

\inferrule{S-IfThen}{
  \eval{E}{e}{true}
}{
  \step{E}{C}{if e then s$_1$ else s$_2$}{E}{C}{s$_1$}%
}

\inferrule{S-IfElse}{
  \eval{E}{e}{false}
}{
  \step{E}{C}{if e then s$_1$ else s$_2$}{E}{C}{s$_2$}%
}

\inferrule{S-WhileDo}{
  \eval{E}{e}{true}
}{
  \step{E}{C}{while e do s}{E}{C}{s; while e do s}
}

\inferrule{S-WhileSkip}{
  \eval{E}{e}{false}
}{
  \step{E}{C}{while e do s}{E}{C}{skip}
}

\section{Babelsberg/Records}
\label{sec:babelsberg-records}


\subsection{Syntax}

The syntax from Section~\ref{babelsberg/int} is augmented now to support
records and the ability to access fields of a record:

\begin{tt}
\begin{tabular}{lrrl}
{\rm Expression} & e & ::= & $\cdots$ $\mid$
\{l$_1$:e$_1$,$\ldots$,l$_n$:e$_n$\} $\mid$ e.l \\
{\rm Label} & l & ::= & {\rm record label names} \\\\
{\rm Value} & v & ::= & c $\mid$ \{l$_1$:v$_1$,$\ldots$,l$_n$:v$_n$\}
\end{tabular}
\end{tt}

We now assume the solver ``understands'' records and record operations
directly.  We assume the equality operator {\tt =} can be used to
compare two records for logical equality, but no other operators apply
to record values.

\subsection{Semantics}


\judgment{\eval{E}{e}{v}}

Expression evaluation is updated to support records. 

\inferrule{E-Rec}{
  \eval{E}{e$_1$}{v$_1$} $\cdots$ \eval{E}{e$_n$}{v$_n$}
}{
  \eval{E}{\{l$_1$:e$_1$,$\ldots$,l$_n$:e$_n$\}}{\{l$_1$:v$_1$,$\ldots$,l$_n$:v$_n$\}}
}

\inferrule{E-Field}{
  \eval{E}{e}{\{l$_1$:v$_1$,$\ldots$,l$_n$:v$_n$\}}
  \andalso 1 $\leq$ $i$ $\leq$ $n$
}{
  \eval{E}{e.l$_i$}{v$_i$}
}


\judgment{\tc{E}{e}{T}}

\judgment{\tcC{E}{C}}

``Expression {\tt e} has type {\tt T} in the context of environment
{\tt E}.''

``Constraint {\tt C} is well formed in the context of environment
{\tt E}.''

We use a notion of typechecking to prevent undesirable nondeterminism
in constraints.  Specifically, we would like for constraint solving to
preserve the structure of the values of variables, changing only the
underlying primitive data as part of a solution.  We formalize our
notion of structure through a simple syntax of types:

\begin{tt}
\begin{tabular}{lrrl}
{\rm Type} & T & ::= & PrimitiveType $\mid$ \{l$_1$:T$_1$,$\ldots$,l$_n$:T$_n$\}
\end{tabular}
\end{tt}

The typechecking rules are mostly standard.  However, we check
expressions dynamically just before constraint solving, so we
typecheck in the context of an runtime environment, which is somewhat unusual.

\inferax{T-Constant}{
  \tc{E}{c}{PrimitiveType}
}

\inferrule{T-Variable}{
  \eval{E}{x}{v} \andalso \tc{E}{v}{T}
}{
  \tc{E}{x}{T}
}

\inferrule{T-ReadOnly}{
  \tc{E}{x}{T}
}{
  \tc{E}{x?}{T}
}

\inferrule{T-Op}{
  \tc{E}{e$_1$}{T} \andalso \tc{E}{e$_2$}{T}
}{
  \tc{E}{e$_1$ $\oplus$ e$_2$}{T}
}

\inferrule{T-Compare}{
  \tc{E}{e$_1$}{T} \andalso \tc{E}{e$_2$}{T}
}{
  \tc{E}{e$_1$ $\olessthan$ e$_2$}{PrimitiveType}
}

The above rule ensures that two records can only be compared for
equality when they have the same structure.  This is necessary to
ensure that the solver will never have to invent record fields in
order to satisfy a record equality constraint.

\inferrule{T-Combine}{
  \tc{E}{e$_1$}{PrimitiveType} \andalso \tc{E}{e$_2$}{PrimitiveType}
}{
  \tc{E}{e$_1$ $\ovee$ e$_2$}{PrimitiveType}
}

\inferrule{T-Rec}{
  \tc{E}{e$_1$}{T$_1$} $\cdots$ \tc{E}{e$_n$}{T$_n$}
}{
  \tc{E}{\{l$_1$:e$_1$,$\ldots$,l$_n$:e$_n$\}}{\{l$_1$:T$_1$,$\ldots$,l$_n$:T$_n$\}}
}

\inferrule{T-Field}{
  \tc{E}{e}{\{l$_1$:T$_1$,$\ldots$,l$_n$:T$_n$\}}
  \andalso $1 \leq i \leq n$
}{
  \tc{E}{e.l$_i$}{T$_i$}
}

The above rule ensures that a field of a record can only be referenced
when it already exists in the record. This is necessary to ensure that
the solver will never have to invent record fields in order to satisfy
a constraint containing field accesses.


The following two rules simply ensure that a constraint is well typed.

\inferrule{T-Constraint}{
  \tc{E}{e}{T}
}{
  {\tcC{E}{e}}
}

\inferrule{T-Priority}{\tc{E}{e}{T}}{\tcC{E}{$\rho$ e}}

\tdm{We should be able to prove that constraint solving preserves
  typing, at least for ``once'' and ``always'' constraints -- not for
  assignments, since there we explicitly allow changing types (see
  below).  And that assignment statements can only change the type of
  the variable that was assigned.  This would be a nice way to validate our choice of typing rules.}

\judgment{\step{E}{C}{s}{E$'$}{C$'$}{s$'$}}

The semantics for executing statements is essentially identical to what we had
before, except that now we typecheck constraints before we solve
them.  Implicitly we get ``stuck'' if such a check fails.
In a practical language, this could be extended to
generate an exception instead.  We only show the modified rules
below.  

\inferrule{S-Asgn}{
  \eval{E}{e}{v} \andalso \tcC{E[x $\mapsto$ v]}{C} \andalso
  \stay{E}{C$_s$} \andalso
  \solve{E$'$}{(C $\wedge$ C$_s$ $\wedge$ x = v)}
}{
  \step{E}{C}{x := e}{E$'$}{C}{skip}
}

The second premise above is necessary so that we don't pass the solver
an ill-typed set of constraints.  Specifically, since the assignment
can change the type of {\tt x}, we have to make sure that the global
constraints in scope, {\tt C}, are well typed with respect to this new
type.  Note that the stay constraint for {\tt x} within {\tt C$_s$} may not be well
typed with respect to the new type of {\tt x}, but that's ok since the
required constraint {\tt x = v} will take precedence.  The syntax {\tt
  E[x $\mapsto$ v]} denotes the environment identical to {\tt E} but
with {\tt x} mapped to the value {\tt v}.

\inferrule{S-AsgnFail}{
  \eval{E}{e}{v} \andalso \tcC{E[x $\mapsto$ v]}{C} \andalso
  \stay{E}{C$_s$} \andalso
  \fail{(C $\wedge$ C$_s$ $\wedge$ x = v)}
}{
  \step{E}{C}{x := e}{E}{C}{fail}
}

\inferrule{S-Once}{\tcC{E}{C$_0$} \andalso
  \stay{E}{C$_s$} \andalso
  \solve{E$'$}{(C $\wedge$ C$_s$ $\wedge$ C$_0$)}
}
{\step{E}{C}{once C$_0$}{E$'$}{C}{skip}}

\inferrule{S-OnceFail}{\tcC{E}{C$_0$} \andalso
  \stay{E}{C$_s$} \andalso
  \fail{(C $\wedge$ C$_s$ $\wedge$ C$_0$)}
}{
  \step{E}{C}{once C$_0$}{E}{C}{fail}
}

\inferrule{S-Always}{\tcC{E}{C$_0$} \andalso
  \stay{E}{C$_s$} \andalso
  \solve{E$'$}{(C $\wedge$ C$_s$ $\wedge$ C$_0$)}
}{
  \step{E}{C}{always C$_0$}{E$'$}{C$'$}{skip}
}

\inferrule{S-AlwaysFail}{\tcC{E}{C$_0$} \andalso
\stay{E}{C$_s$} \andalso
  \fail{(C $\wedge$ C$_s$ $\wedge$ C$_0$)}
}{
  \step{E}{C}{always C$_0$}{E}{C}{fail}
}




\end{document}

